% =====================================================================
%  3.4 Data splitting
% =====================================================================
\subsection{Data Splitting}
\label{sec:splitting}

\begin{table}[H]
  \centering
  \small
  \caption{Chronological split. Random splits would leak future information into training.}
  \label{tab:split}
  \begin{tabularx}{\linewidth}{@{}l l r Y@{}}
    \toprule
    Set & Period & Usable hours & Why \\
    \midrule
    Train & 2015-01-08 -- 2016-12-31 & 17{,}332 & Two winters and two summers; the first week has no
      $y_{D-7,h}$. \\
    Validation & 2017 & 8{,}757 & A full year, so every season counts when choosing
      hyper-parameters. \\
    Test & 2018 & 8{,}756 & Touched once, after the configuration is frozen. \\
    \bottomrule
  \end{tabularx}
\end{table}

All hyper-parameters are chosen on validation with models fitted on 2015--2016. For the test, each
model is refitted \emph{once} on 2015--2017 with the frozen hyper-parameters and predicts all of
2018. This repeats the validation protocol (fit on earlier years, forecast the next full year), so
validation and test errors are comparable. Both weather settings use exactly the same rows, so all
models are scored on the same hours.

Two safeguards protect the test year. First, the development scripts never load 2018. Second, the
test script asserts that adding 2018 changes no feature of any earlier row, and it refuses to run
if the frozen parameters come from a smoke run or if the test has already been run.
